\srcpage{323}
\section*{第14章\quad 拟阵的推广}\addcontentsline{toc}{section}{第14章\quad 拟阵的推广}
拟阵有若干有趣的推广。第13.1节已经见到独立系统，它来自删去公理(M3)。第14.1节考虑贪心结构，它来自改为删去(M2)。此外，与拟阵和次模函数有关的一类多胞体，称为多拟阵，可将重要定理大幅推广；第14.2节将讨论它们。第14.3、14.4节考虑任意次模函数最小化问题的两种方法：一种使用椭球法，另一种使用组合算法。对于对称次模函数这一重要特例，第14.5节介绍一种更简单的算法。
\sect{14.1\quad 贪心结构}
由定义，集合系统 $(E,\F)$ 是拟阵，当且仅当满足：
\begin{description}
\item[(M1)] $\varnothing\in\F$；
\item[(M2)] 若 $X\subseteq Y\in\F$，则 $X\in\F$；
\item[(M3)] 若 $X,Y\in\F$ 且 $|X|>|Y|$，则存在 $x\in X\setminus Y$，使 $Y\cup\{x\}\in\F$。
\end{description}
删去(M3)就得到第13.1、13.4节讨论的独立系统。现在改为删去(M2)。
\st{定义14.1} \textbf{贪心结构}（greedoid）是满足(M1)、(M3)的集合系统 $(E,\F)$。

取代子集封闭性(M2)的是可达性：若 $\varnothing\in\F$，且对每个 $X\in\F\setminus\{\varnothing\}$，存在 $x\in X$ 使 $X\setminus\{x\}\in\F$，则称 $(E,\F)$ \textbf{可达}。贪心结构具有可达性（由(M1)、(M3)直接得到）。虽然比拟阵更一般，但它们仍有丰富的结构，另一方面也推广了许多不同的、表面上没有联系的概念。从以下结果开始。
\st{定理14.2} 设 $(E,\F)$ 是可达集合系统。以下命题等价：
\begin{enumerate}[label=(\alph*)]
\item 对任意 $X\subseteq Y\subset E$ 及 $z\in E\setminus Y$，若 $X\cup\{z\}\in\F$ 且 $Y\in\F$，则 $Y\cup\{z\}\in\F$；
\item $\F$ 对并运算封闭。
\end{enumerate}

\srcpage{324}
\pr (a)$\Rightarrow$(b)：设 $X,Y\in\F$，证明 $X\cup Y\in\F$。取满足 $Z\in\F$、$X\subseteq Z\subseteq X\cup Y$ 的极大集合 $Z$。假设 $Y\setminus Z\ne\varnothing$。对 $Y$ 反复应用可达性，可得 $Y'\in\F$、$Y'\subseteq Z$，以及 $y\in Y\setminus Z$，使 $Y'\cup\{y\}\in\F$。对 $Z,Y',y$ 应用(a)，得 $Z\cup\{y\}\in\F$，与 $Z$ 的选取矛盾。(b)$\Rightarrow$(a)显然。\qed
若满足定理14.2的条件，则称 $(E,\F)$ 为\textbf{反拟阵}（antimatroid）。
\st{命题14.3} 每个反拟阵都是贪心结构。
\pr 设 $(E,\F)$ 是反拟阵，即可达且对并封闭。为证(M3)，设 $X,Y\in\F$、$|X|>|Y|$。由可达性，可将 $X$ 排为 $\{x_1,\ldots,x_n\}$，使 $\{x_1,\ldots,x_i\}\in\F$（$i=0,\ldots,n$）。取最小的 $i$ 使 $x_i\notin Y$。因 $\F$ 对并封闭，$Y\cup\{x_i\}=Y\cup\{x_1,\ldots,x_i\}\in\F$。\qed
反拟阵也可用闭包算子等价定义。
\st{命题14.4} 设集合系统 $(E,\F)$ 满足 $\F$ 对并封闭且 $\varnothing\in\F$。定义
\[\tau(A):=\bigcap\{X\subseteq E:A\subseteq X,\ E\setminus X\in\F\}.\]
则 $\tau$ 是闭包算子，即满足定理13.11的(S1)--(S3)。
\pr 设 $X\subseteq Y\subseteq E$。$X\subseteq\tau(X)\subseteq\tau(Y)$ 显然。为证(S3)，假设存在 $y\in\tau(\tau(X))\setminus\tau(X)$。则对所有满足 $\tau(X)\subseteq Y$、$E\setminus Y\in\F$ 的 $Y\subseteq E$，都有 $y\in Y$；但存在 $Z\subseteq E\setminus\{y\}$，使 $X\subseteq Z$、$E\setminus Z\in\F$。这蕴含 $\tau(X)\nsubseteq Z$，矛盾。\qed
\st{定理14.5} 设集合系统 $(E,\F)$ 满足 $\F$ 对并封闭且 $\varnothing\in\F$。则 $(E,\F)$ 可达，当且仅当命题14.4的闭包算子 $\tau$ 满足\textbf{反交换性质}：若 $X\subseteq E$、$y,z\in E\setminus\tau(X)$、$y\ne z$，且 $z\in\tau(X\cup\{y\})$，则 $y\notin\tau(X\cup\{z\})$。
\pr 若 $(E,\F)$ 可达，则由命题14.3，(M3)成立。为证反交换性质，设 $X\subseteq E$、$B:=E\setminus\tau(X)$，以及 $y,z\in B$，满足 $z\notin A:=E\setminus\tau(X\cup\{y\})$。注意 $A,B\in\F$，且 $A\subseteq B\setminus\{y,z\}$。对 $A,B$ 应用(M3)，得 $b\in B\setminus A\subseteq E\setminus(X\cup A)$，使 $A\cup\{b\}\in\F$。$A\cup\{b\}$ 不可能包含在 $E\setminus(X\cup\{y\})$ 中，否则 $\tau(X\cup\{y\})\subseteq E\setminus(A\cup\{b\})$，与 $\tau(X\cup\{y\})=E\setminus A$ 矛盾。因此 $b=y$。所以 $A\cup\{y\}\in\F$，从而 $\tau(X\cup\{z\})\subseteq E\setminus(A\cup\{y\})$，得到 $y\notin\tau(X\cup\{z\})$。

反之，取 $A\in\F\setminus\{\varnothing\}$，令 $X:=E\setminus A$，则 $\tau(X)=X$。取 $a\in A$ 使 $|\tau(X\cup\{a\})|$ 最小。断言 $\tau(X\cup\{a\})=X\cup\{a\}$，即 $A\setminus\{a\}\in\F$。否则存在 $b\in\tau(X\cup\{a\})\setminus(X\cup\{a\})$。由反交换性质，$a\notin\tau(X\cup\{b\})$。此外，

\srcpage{325}
\[\tau(X\cup\{b\})\subseteq\tau(\tau(X\cup\{a\})\cup\{b\})
=\tau(\tau(X\cup\{a\}))=\tau(X\cup\{a\}).\]
故 $\tau(X\cup\{b\})$ 是 $\tau(X\cup\{a\})$ 的真子集，与 $a$ 的选取矛盾。\qed
定理14.5的反交换性质与(S4)不同。定理13.11的(S4)是 $\R^n$ 中线性包的一条性质，而反交换性质是 $\R^n$ 中凸包的一条性质：若 $y\ne z$、$z\notin\conv(X)$ 且 $z\in\conv(X\cup\{y\})$，则显然 $y\notin\conv(X\cup\{z\})$。所以对任意有限集 $E\subset\R^n$，
\[(E,\{X\subseteq E:X\cap\conv(E\setminus X)=\varnothing\})\]
是反拟阵。

贪心结构推广了拟阵和反拟阵，还包含其他有趣结构。例如 Edmonds 基数匹配算法中使用的花结构（习题1）。另一个基本例子是：
\st{命题14.6} 设 $G$ 是有向或无向图，$r\in V(G)$。令 $\F$ 是 $G$ 中所有以 $r$ 为根的树形图的边集族，或所有包含 $r$ 的树的边集族（不要求生成）。则 $(E(G),\F)$ 是贪心结构。
\pr (M1)显然。对有向情形证明(M3)，无向情形同理。设 $(X_1,F_1)$、$(X_2,F_2)$ 是以 $r$ 为根的两个树形图，且 $|F_1|>|F_2|$。则 $|X_1|=|F_1|+1>|F_2|+1=|X_2|$，取 $x\in X_1\setminus X_2$。$(X_1,F_1)$ 中的 $r$--$x$ 路包含某条边 $(v,w)$，其中 $v\in X_2$、$w\notin X_2$。可将该边加入 $(X_2,F_2)$，所以 $F_2\cup\{(v,w)\}\in\F$。\qed
这称为 $G$ 的\textbf{有向（无向）分支贪心结构}。

在具有非负权的连通图 $G$ 中寻找最大权生成树，是圈拟阵 $M(G)$ 上的最大化问题。此时择优加入贪心算法就是 Kruskal 算法。现在同一问题有了第二种表述：在 $G$ 的无向分支贪心结构 $(E(G),\F)$ 中，寻找最大权集合 $F\in\F$。

下面给出贪心结构上的一般贪心算法。在拟阵特例下，它恰是第13.4节的择优加入贪心算法；对具有模费用函数 $c$ 的无向分支贪心结构，它就是 Prim 算法。
\alg{贪心结构的贪心算法}
\textbf{输入：}贪心结构 $(E,\F)$ 及函数 $c:2^E\to\R$，由预言机给出：对给定 $X\subseteq E$，判定 $X\in\F$ 并返回 $c(X)$。\par
\textbf{输出：}集合 $F\in\F$。
\begin{enumerate}[label=\arabic*.]
\item 令 $F:=\varnothing$。
\item 选 $e\in E\setminus F$，使 $F\cup\{e\}\in\F$ 且 $c(F\cup\{e\})$ 最大。若不存在则停止。
\item 令 $F:=F\cup\{e\}$，返回步骤2。
\end{enumerate}

\srcpage{326}
即使 $c$ 是模费用函数，该算法也不总给出最优解。至少可以刻画它在哪些贪心结构上有效。
\st{定理14.7} 设 $(E,\F)$ 是贪心结构。贪心结构的贪心算法对每个模权函数 $c:2^E\to\R_+$ 都能找到最大权集合 $F\in\F$，当且仅当 $(E,\F)$ 具有所谓的\textbf{强交换性质}：对所有 $A\in\F$、$\F$ 中极大的 $B$、$A\subseteq B$，以及满足 $A\cup\{x\}\in\F$ 的 $x\in E\setminus B$，存在 $y\in B\setminus A$，使 $A\cup\{y\}\in\F$ 且 $(B\setminus\{y\})\cup\{x\}\in\F$。
\pr 设 $(E,\F)$ 具有强交换性质。取 $c:E\to\R_+$，设算法所得解为 $A=\{a_1,\ldots,a_l\}$，元素按 $a_1,\ldots,a_l$ 的顺序选入。取最优解 $B=\{a_1,\ldots,a_k\}\mathbin{\dot\cup}B'$，使 $k$ 最大，假设 $k<l$。对 $\{a_1,\ldots,a_k\}$、$B$、$a_{k+1}$ 应用强交换性质，得 $y\in B'$，使 $\{a_1,\ldots,a_k,y\}\in\F$ 且 $(B\setminus\{y\})\cup\{a_{k+1}\}\in\F$。由算法步骤2对 $a_{k+1}$ 的选择，$c(a_{k+1})\ge c(y)$，从而 $c((B\setminus\{y\})\cup\{a_{k+1}\})\ge c(B)$，与 $B$ 的选取矛盾。

反之，设 $(E,\F)$ 不具有强交换性质。取 $A\in\F$、$\F$ 中极大的 $B$、$A\subseteq B$ 和 $x\in E\setminus B$，满足 $A\cup\{x\}\in\F$，但所有满足 $A\cup\{y\}\in\F$ 的 $y\in B\setminus A$ 都有 $(B\setminus\{y\})\cup\{x\}\notin\F$。
令 $Y:=\{y\in B\setminus A:A\cup\{y\}\in\F\}$。定义 $c(e):=2$（$e\in B\setminus Y$）、$c(e):=1$（$e\in Y\cup\{x\}$）、$c(e):=0$（$e\in E\setminus(B\cup\{x\})$）。算法可能先选 $A$ 的元素（权重为2），再选 $x$。最终所得 $F\in\F$ 不可能最优，因为
\[c(F)\le c(B\cup\{x\})-2<c(B\cup\{x\})-1=c(B),\qquad B\in\F.\]\qed
事实上，在一般贪心结构上优化模函数是 NP 困难的。这可由以下观察结合推论15.24得到。
\st{命题14.8} 给定无向图 $G$、$k\in\N$，判定 $G$ 是否有基数为 $k$ 的顶点覆盖，这个问题可线性归约为：给定由隶属预言机描述的贪心结构 $(E,\F)$ 及函数 $c:E\to\R_+$，求使 $c(F)$ 最大的 $F\in\F$。
\pr 任取无向图 $G$、$k\in\N$。令 $D:=V(G)\mathbin{\dot\cup}E(G)$，并令
\[\F:=\{X\subseteq D:\text{对所有 }e=\{v,w\}\in E(G)\cap X,\ v\in X\text{ 或 }w\in X\}.\]
$(D,\F)$ 可达且对并封闭，因此是反拟阵，特别由命题14.3，它是贪心结构。
考虑 $\F':=\{X\in\F:|X|\le|E(G)|+k\}$。因(M1)、(M3)保留，$(D,\F')$ 仍是贪心结构。令 $c(e):=1$（$e\in E(G)$）、$c(v):=0$（$v\in V(G)$）。则存在 $F\in\F'$ 满足 $c(F)=|E(G)|$，当且仅当 $G$ 有大小为 $k$ 的顶点覆盖。\qed

\srcpage{327}
另一方面，有些有趣函数能在任意贪心结构上最大化，例如对 $c':E\to\R_+$，瓶颈函数 $c(F):=\min\{c'(e):e\in F\}$（习题2）。更多结果见 Korte、Lovász 和 Schrader [1991]。
\sect{14.2\quad 多拟阵}
由定理13.10可知拟阵与次模函数的紧密联系。次模函数定义了以下有趣的多面体类。
\st{定义14.9} \textbf{多拟阵}（polymatroid）是如下形式的多胞体：
\[P(f):=\left\{x\in\R^E:x\ge0,\ \sum_{e\in A}x_e\le f(A)\text{ 对所有 }A\subseteq E\right\},\]
其中 $E$ 为有限集，$f:2^E\to\R_+$ 是次模函数。
不难看出，对任意多拟阵，都可选取满足 $f(\varnothing)=0$ 且单调的 $f$（习题5；若 $X\subseteq Y\subseteq E$ 蕴含 $f(X)\le f(Y)$，则称 $f:2^E\to\R$ \textbf{单调}）。Edmonds 的原定义有所不同，见习题6。还需指出，polymatroid 一词有时用于二元组 $(E,f)$，而不是该多胞体。

若 $f$ 是某拟阵的秩函数，则 $P(f)$ 是其独立集关联向量的凸包（定理13.21）。已知择优加入贪心算法能在拟阵多胞体上优化任意线性函数。类似的贪心算法也适用于一般多拟阵。假定 $f$ 单调。
\alg{多拟阵贪心算法}
\textbf{输入：}有限集 $E$，由预言机给出的次模、单调函数 $f:2^E\to\R_+$，以及向量 $c\in\R^E$。\par
\textbf{输出：}使 $cx$ 最大的向量 $x\in P(f)$。
\begin{enumerate}[label=\arabic*.]
\item 将 $E=\{e_1,\ldots,e_n\}$ 排序，使 $c(e_1)\ge\cdots\ge c(e_k)>0\ge c(e_{k+1})\ge\cdots\ge c(e_n)$。
\item 若 $k\ge1$，令 $x(e_1):=f(\{e_1\})$；对 $i=2,\ldots,k$，令
\[x(e_i):=f(\{e_1,\ldots,e_i\})-f(\{e_1,\ldots,e_{i-1}\});\]
对 $i=k+1,\ldots,n$，令 $x(e_i):=0$。
\end{enumerate}
\st{命题14.10} 设 $E=\{e_1,\ldots,e_n\}$，$f:2^E\to\R$ 是次模函数，且 $f(\varnothing)\ge0$。设 $b:E\to\R$，满足 $b(e_1)\le f(\{e_1\})$，且对 $i=2,\ldots,n$，$b(e_i)\le f(\{e_1,\ldots,e_i\})-f(\{e_1,\ldots,e_{i-1}\})$。则对所有 $A\subseteq E$，有 $\sum_{a\in A}b(a)\le f(A)$。

\srcpage{328}
\pr 对 $i=\max\{j:e_j\in A\}$ 归纳。$A=\varnothing$ 或 $A=\{e_1\}$ 时显然。若 $i\ge2$，则
\begin{align*}
\sum_{a\in A}b(a)&=\sum_{a\in A\setminus\{e_i\}}b(a)+b(e_i)\le f(A\setminus\{e_i\})+b(e_i)\\
&\le f(A\setminus\{e_i\})+f(\{e_1,\ldots,e_i\})-f(\{e_1,\ldots,e_{i-1}\})\le f(A),
\end{align*}
第一个不等式来自归纳假设，第三个来自次模性。\qed
\st{定理14.11} 多拟阵贪心算法能正确求得使 $cx$ 最大的 $x\in P(f)$。若 $f$ 为整数值函数，则 $x$ 也是整数向量。
\pr 设 $x\in\R^E$ 是算法对 $E,f,c$ 的输出。由定义，若 $f$ 为整数值，则 $x$ 也为整数。因为 $f$ 单调，有 $x\ge0$，故由命题14.10，$x\in P(f)$。
现设 $y\in\R_+^E$ 且 $cy>cx$。类似于定理13.19的证明，令 $d_j:=c(e_j)-c(e_{j+1})$（$j=1,\ldots,k-1$）、$d_k:=c(e_k)$，有
\[\sum_{j=1}^kd_j\sum_{i=1}^jx(e_i)=cx<cy\le\sum_{j=1}^kc(e_j)y(e_j)
=\sum_{j=1}^kd_j\sum_{i=1}^jy(e_i).\]
因所有 $d_j\ge0$，存在 $j\in\{1,\ldots,k\}$，使 $\sum_{i=1}^jy(e_i)>\sum_{i=1}^jx(e_i)$。但 $\sum_{i=1}^jx(e_i)=f(\{e_1,\ldots,e_j\})$，故 $y\notin P(f)$。\qed
与拟阵一样，也能处理两个多拟阵的交。下面的多拟阵交定理有许多推论。
\st{定理14.12（Edmonds [1970,1979]）} 设 $E$ 是有限集，$f,g:2^E\to\R_+$ 是次模函数。则下述系统是全对偶整数的（TDI）：
\[x\ge0,\qquad\sum_{e\in A}x_e\le f(A)\ (A\subseteq E),\qquad
\sum_{e\in A}x_e\le g(A)\ (A\subseteq E).\]
\pr 考虑下列互为原始与对偶的线性规划：
\[\max\left\{cx:\sum_{e\in A}x_e\le f(A),\ \sum_{e\in A}x_e\le g(A)\text{ 对所有 }A\subseteq E,\ x\ge0\right\},\]
以及
\[\min\left\{\sum_{A\subseteq E}(f(A)y_A+g(A)z_A):\sum_{\substack{A\subseteq E\\e\in A}}(y_A+z_A)\ge c_e\text{ 对所有 }e\in E,\ y,z\ge0\right\}.\]
为证明全对偶整数性，使用引理5.22。

\srcpage{329}
设 $c:E\to\Z$，取最优对偶解 $y,z$，使
\[\sum_{A\subseteq E}(y_A+z_A)|A|\,|E\setminus A|\tag{14.1}\]
尽可能小。断言 $\F:=\{A\subseteq E:y_A>0\}$ 是链，即任意 $A,B\in\F$ 都满足 $A\subseteq B$ 或 $B\subseteq A$。
否则取 $A,B\in\F$，使 $A\cap B\ne A,B$。令 $\epsilon:=\min\{y_A,y_B\}$，令 $y'_A:=y_A-\epsilon$、$y'_B:=y_B-\epsilon$、$y'_{A\cap B}:=y_{A\cap B}+\epsilon$、$y'_{A\cup B}:=y_{A\cup B}+\epsilon$，其他分量不变。$y',z$ 是对偶可行解；由 $f$ 的次模性，它也最优。但对它而言(14.1)更小，与 $y$ 的选取矛盾。

同理 $\F':=\{A\subseteq E:z_A>0\}$ 也是链。令 $M,M'$ 为列由 $E$ 的元素标号、行分别为 $\F,\F'$ 中集合的关联向量的矩阵。由引理5.22，只需证明 $\binom M{M'}$ 全幺模。
使用 Ghouila-Houri 定理5.23。任取一组行 $\mathcal R=\{A_1,\ldots,A_p,B_1,\ldots,B_q\}$，其中 $A_1\supseteq\cdots\supseteq A_p$、$B_1\supseteq\cdots\supseteq B_q$。令
\[\mathcal R_1:=\{A_i:i\text{ 为奇数}\}\cup\{B_i:i\text{ 为偶数}\},\qquad
\mathcal R_2:=\mathcal R\setminus\mathcal R_1.\]
对任意 $e\in E$，包含它的行构成 $\{A_1,\ldots,A_{p_e}\}\cup\{B_1,\ldots,B_{q_e}\}$，其中 $p_e\in\{0,\ldots,p\}$、$q_e\in\{0,\ldots,q\}$。因此 $\mathcal R_1$ 中各行之和减去 $\mathcal R_2$ 中各行之和，所得向量只含 $-1,0,1$。定理5.23的判据成立。\qed
可以在两个多拟阵之交上优化线性函数，不过不像单个多拟阵那样容易。如果能解决每个多拟阵的分离问题，就可使用椭球法。第14.3节将再次讨论。
\st{推论14.13（Edmonds [1970]）} 设 $(E,\mathcal M_1)$、$(E,\mathcal M_2)$ 是秩函数分别为 $r_1,r_2$ 的两个拟阵。则 $\mathcal M_1\cap\mathcal M_2$ 各元素关联向量的凸包为
\[\left\{x\in\R_+^E:\sum_{e\in A}x_e\le\min\{r_1(A),r_2(A)\}\text{ 对所有 }A\subseteq E\right\}.\]
\pr $r_1,r_2$ 非负且次模（定理13.10），故上述不等式系统 TDI（定理14.12）。因为 $r_1,r_2$ 是整数值函数，多胞体为整数多胞体（推论5.14）。由于对所有 $A\subseteq E$ 都有 $r_1(A)\le|A|$，其顶点均为0--1向量，因而是公共独立集（$\mathcal M_1\cap\mathcal M_2$ 中元素）的关联向量。由推论3.27，多胞体就是这些顶点的凸包。反之，每个这样的关联向量都满足不等式（由秩函数定义）。\qed
取 $\mathcal M_1=\mathcal M_2$，当然得到拟阵多胞体的描述（定理13.21）。定理14.12还有以下推论。

\srcpage{330}
\st{推论14.14（Edmonds [1970]）} 设 $E$ 是有限集，$f,g:2^E\to\R_+$ 次模且单调。则
\[\max\{\one x:x\in P(f)\cap P(g)\}=\min_{A\subseteq E}(f(A)+g(E\setminus A)).\]
此外，若 $f,g$ 为整数值函数，则存在达到最大值的整数向量 $x$。
\pr 由定理14.12，$\max\{\one x:x\in P(f)\cap P(g)\}$ 的对偶问题，即
\[\min\left\{\sum_{A\subseteq E}(f(A)y_A+g(A)z_A):\sum_{\substack{A\subseteq E\\e\in A}}(y_A+z_A)\ge1\text{ 对所有 }e\in E,\ y,z\ge0\right\},\]
具有整数最优解 $y,z$。令 $B:=\bigcup_{A:y_A\ge1}A$、$C:=\bigcup_{A:z_A\ge1}A$。令 $y'_B:=1$、$z'_C:=1$，其余分量均为0。则 $B\cup C=E$，且 $y',z'$ 是对偶可行解。由 $f,g$ 非负且次模，
\[\sum_{A\subseteq E}(f(A)y_A+g(A)z_A)\ge f(B)+g(C).\]
因为 $E\setminus B\subseteq C$ 且 $g$ 单调，右端至少为 $f(B)+g(E\setminus B)$，故“$\ge$”得证。
另一方向“$\le$”显然：任取 $A\subseteq E$，令 $y_A:=1$、$z_{E\setminus A}:=1$，其他分量均为0，即得对偶可行解。整数性直接来自定理14.12与推论5.14。\qed
定理13.31是其特例。此外还得到：
\st{推论14.15（Frank [1982]）} 设 $E$ 是有限集，$f,g:2^E\to\R$，其中 $f$ 超模、$g$ 次模，且 $f\le g$。则存在模函数 $h:2^E\to\R$，使 $f\le h\le g$。若 $f,g$ 为整数值函数，则可取整数值的 $h$。
\pr 令 $M:=2\max\{|f(A)|+|g(A)|:A\subseteq E\}$，并对所有 $A\subseteq E$ 定义
\[f'(A):=g(E)-f(E\setminus A)+M|A|,\qquad g'(A):=g(A)-f(\varnothing)+M|A|.\]
$f',g'$ 非负、次模且单调。应用推论14.14，得
\begin{align*}
\max\{\one x:x\in P(f')\cap P(g')\}
&=\min_{A\subseteq E}(f'(A)+g'(E\setminus A))\\
&=\min_{A\subseteq E}\bigl(g(E)-f(E\setminus A)+M|A|\\
&\qquad+g(E\setminus A)-f(\varnothing)+M|E\setminus A|\bigr)\\
&\ge g(E)-f(\varnothing)+M|E|.
\end{align*}
因此取 $x\in P(f')\cap P(g')$，使 $\one x=g(E)-f(\varnothing)+M|E|$。若 $f,g$ 为整数值，可取整数 $x$。令 $h'(A):=\sum_{e\in A}x_e$，并对所有 $A\subseteq E$ 令 $h(A):=h'(A)+f(\varnothing)-M|A|$。
\srcpage{331}
$h$ 是模函数。此外，对所有 $A\subseteq E$，有
\[h(A)\le g'(A)+f(\varnothing)-M|A|=g(A),\]
\begin{align*}
h(A)&=\one x-h'(E\setminus A)+f(\varnothing)-M|A|\\
&\ge g(E)+M|E|-M|A|-f'(E\setminus A)=f(A).
\end{align*}\qed
这与凸函数和凹函数的类比显而易见；另见习题9。
\sect{14.3\quad 次模函数最小化}
多拟阵 $P(f)$ 对向量 $x$ 的分离问题要求找出集合 $A$，使 $f(A)<\sum_{e\in A}x(e)$。故可归约为寻找使 $g(A)$ 最小的集合 $A$，其中 $g(A):=f(A)-\sum_{e\in A}x(e)$。注意若 $f$ 次模，则 $g$ 也次模。因此次模函数最小化是一个有趣的问题。

另一动机是次模函数可视为凸函数的离散对应物（推论14.15及习题9）。第8.7节已经解决了一个特例：在无向图中寻找最小割，可视为在 $2^U\setminus\{\varnothing,U\}$ 上最小化某个对称次模函数 $f:2^U\to\R_+$。回到该特例之前，先说明如何最小化一般次模函数。假定已给出 $\operatorname{size}(f(S))$ 的上界。为简单起见，仅考虑整数值次模函数。
\alg{次模函数最小化问题}
\textbf{实例：}有限集 $U$，由预言机给出的次模函数 $f:2^U\to\Z$。\par
\textbf{任务：}求使 $f(X)$ 最小的子集 $X\subseteq U$。

Grötschel、Lovász 和 Schrijver [1981] 说明了如何借助椭球法求解。其思想是通过二分搜索确定最小值，从而将问题归约为某个多拟阵的分离问题。利用分离与优化的等价性（第4.6节），只需在多拟阵上优化线性函数；而多拟阵贪心算法很容易做到这一点。首先需要 $S\subseteq U$ 时 $|f(S)|$ 的上界。
\st{命题14.16} 对任意次模函数 $f:2^U\to\Z$ 及任意 $S\subseteq U$，
\[f(U)-\sum_{u\in U}\max\{0,f(\{u\})-f(\varnothing)\}
\le f(S)\le f(\varnothing)+\sum_{u\in U}\max\{0,f(\{u\})-f(\varnothing)\}.\]
特别地，只需 $|U|+2$ 次对 $f$ 的预言机调用，就能在线性时间内计算一个数 $B$，使所有 $S\subseteq U$ 都有 $|f(S)|\le B$。
\pr 对 $\varnothing\ne S\subseteq U$，取 $x\in S$，反复应用次模性可得：

\srcpage{332}
\[f(S)\le-f(\varnothing)+f(S\setminus\{x\})+f(\{x\})
\le\cdots\le-|S|f(\varnothing)+f(\varnothing)+\sum_{x\in S}f(\{x\}).\]
对 $S\subset U$，取 $y\in U\setminus S$，有
\begin{align*}
f(S)&\ge-f(\{y\})+f(S\cup\{y\})+f(\varnothing)\\
&\ge\cdots\ge-\sum_{y\in U\setminus S}f(\{y\})+f(U)+|U\setminus S|f(\varnothing).
\end{align*}\qed
\st{命题14.17} 以下问题可在多项式时间内求解：给定有限集 $U$，由预言机给出的次模、单调函数 $f:2^U\to\Z_+$，且 $S\ne\varnothing$ 时 $f(S)>0$；给定 $B\in\N$，使所有 $S\subseteq U$ 满足 $f(S)\le B$；再给定 $x\in\Z_+^U$。判定 $x\in P(f)$，若不成立则返回一个 $S\subseteq U$，使 $\sum_{v\in S}x(v)>f(S)$。
\pr 这就是多拟阵 $P(f)$ 的分离问题。使用定理4.23，因为已解决 $P(f)$ 上的优化问题：多拟阵贪心算法能最大化其上的任意线性函数（定理14.11）。

需检验定理4.23的前提。因零向量及所有单位向量都在 $P(f)$ 中，可取内点 $x_0:=\epsilon\one$，其中 $\epsilon=1/(|U|+1)$。有 $\operatorname{size}(x_0)=O(|U|\log|U|)$。此外，$P(f)$ 的每个顶点都可由多拟阵贪心算法对某个目标函数产生（参见定理14.11），故其编码长度为 $O(|U|(2+\log B))$。因此分离问题可在多项式时间内解决。由定理4.23，若 $x\notin P(f)$，就得到一个被 $x$ 违反的、定义 $P(f)$ 某个面的不等式，对应于某个满足 $\sum_{v\in S}x(v)>f(S)$ 的 $S\subseteq U$。\qed
由于我们不要求 $f$ 单调，不能直接使用此结果，而考虑另一函数。
\st{命题14.18} 设 $f:2^U\to\R$ 次模，$\beta\in\R$。定义 $g:2^U\to\R$ 为
\[g(X):=f(X)-\beta+\sum_{e\in X}(f(U\setminus\{e\})-f(U)),\]
则 $g$ 次模且单调。
\pr 次模性直接来自 $f$ 的次模性。为证单调性，设 $X\subset U$、$e\in U\setminus X$。由次模性，
\[g(X\cup\{e\})-g(X)=f(X\cup\{e\})-f(X)+f(U\setminus\{e\})-f(U)\ge0.\]\qed
\st{定理14.19} 次模函数最小化问题可在关于 $|U|+\log\max\{|f(S)|:S\subseteq U\}$ 的多项式时间内求解。
\pr 设 $U$ 有限，$f$ 由预言机给出。先计算 $B\in\N$，使所有 $S\subseteq U$ 都满足 $|f(S)|\le B$（命题14.16）。由 $f$ 的次模性，对每个 $e\in U$、$X\subseteq U\setminus\{e\}$，有：

\srcpage{333}
\[f(\{e\})-f(\varnothing)\ge f(X\cup\{e\})-f(X)\ge f(U)-f(U\setminus\{e\}).\tag{14.2}\]
若某个 $e\in U$ 满足 $f(\{e\})-f(\varnothing)\le0$，则由(14.2)，存在包含 $e$ 的最优集合 $S$。此时考虑实例 $(U',B,f')$，其中 $U':=U\setminus\{e\}$，$f'(X):=f(X\cup\{e\})$（$X\subseteq U\setminus\{e\}$）。求出使 $f'(S')$ 最小的 $S'\subseteq U'$，并输出 $S:=S'\cup\{e\}$。

同样，若 $f(U)-f(U\setminus\{e\})\ge0$，则由(14.2)，存在不含 $e$ 的最优集合 $S$。此时只需在 $U\setminus\{e\}$ 上最小化 $f$。两种情形都缩小了底集。

故可假定所有 $e\in U$ 都满足 $f(\{e\})-f(\varnothing)>0$、$f(U\setminus\{e\})-f(U)>0$。令 $x(e):=f(U\setminus\{e\})-f(U)$。对每个整数 $-B\le\beta\le f(\varnothing)$，定义 $g(X):=f(X)-\beta+\sum_{e\in X}x(e)$。由命题14.18，$g$ 次模且单调。又有 $g(\varnothing)=f(\varnothing)-\beta\ge0$，且所有 $e\in U$ 满足 $g(\{e\})=f(\{e\})-\beta+x(e)>0$，故非空 $X$ 均有 $g(X)>0$。应用命题14.17检验 $x\in P(g)$。若成立，则所有 $X\subseteq U$ 都满足 $f(X)\ge\beta$，此次判定完成；否则得到满足 $f(S)<\beta$ 的集合 $S$。

现作二分搜索，每次适当选择 $\beta$，经 $O(\log(2B))$ 次迭代，找到 $\beta^*\in\{-B,-B+1,\ldots,f(\varnothing)\}$，使所有 $X\subseteq U$ 满足 $f(X)\ge\beta^*$，但某个 $S\subseteq U$ 满足 $f(S)<\beta^*+1$。此 $S$ 即最小化 $f$。\qed
首个强多项式时间算法由 Grötschel、Lovász 和 Schrijver [1988] 设计，同样基于椭球法。Schrijver [2000] 以及独立工作的 Iwata、Fleischer 和 Fujishige [2001] 找到了强多项式时间的组合算法。下一节描述 Schrijver 算法。
\sect{14.4\quad Schrijver 算法}
给定有限集 $U$ 及次模函数 $f:2^U\to\R$，不妨设 $U=\{1,\ldots,n\}$ 且 $f(\varnothing)=0$。Schrijver [2000] 的算法始终维持 $f$ 的所谓\textbf{基多面体}中的一个点 $x$。该多面体定义为
\[\left\{x\in\R^U:\sum_{u\in A}x(u)\le f(A)\text{ 对所有 }A\subseteq U,\ \sum_{u\in U}x(u)=f(U)\right\}.\]
其顶点集合恰为 $U$ 的所有全序 $\prec$ 对应的向量 $b^\prec$，其中
\[b^\prec(u):=f(\{v\in U:v\preceq u\})-f(\{v\in U:v\prec u\})\quad(u\in U).\]
这里不会用到这一事实；它可类似于定理14.11那样证明（习题13）。

\srcpage{334}
点 $x$ 始终显式写成这些顶点的凸组合 $x=\lambda_1b^{\prec_1}+\cdots+\lambda_kb^{\prec_k}$。初始可取 $k=1$ 及任意全序。
\alg{Schrijver 算法}
\textbf{输入：}有限集 $U=\{1,\ldots,n\}$；由预言机给出的次模函数 $f:2^U\to\Z$，满足 $f(\varnothing)=0$。\par
\textbf{输出：}使 $f(X)$ 最小的子集 $X\subseteq U$。
\begin{enumerate}[label=\arabic*.]
\item 令 $k:=1$，取 $U$ 上任意全序 $\prec_1$，令 $x:=b^{\prec_1}$。
\item 令 $D:=(U,A)$，其中 $A=\{(u,v):\text{某个 }i\in\{1,\ldots,k\}\text{ 满足 }u\prec_i v\}$。
\item 令 $P:=\{v\in U:x(v)>0\}$、$N:=\{v\in U:x(v)<0\}$，令 $X$ 为 $D$ 中从 $P$ 不可达的顶点集。若 $N\subseteq X$，则停止。否则以 $d(v)$ 表示 $D$ 中 $P$ 到 $v$ 的距离。
\item 在从 $P$ 可达的 $t\in N$ 中，选择使 $(d(t),t)$ 按字典序最大的顶点；再选择满足 $(s,t)\in A$、$d(s)=d(t)-1$ 且 $s$ 最大的顶点 $s$。选 $i\in\{1,\ldots,k\}$，使 $\alpha:=|\{v:s\prec_i v\preceq_i t\}|$ 最大（达到此最大值的下标数目记为 $\beta$）。
\item 以 $\prec_i^{s,u}$ 表示将 $u$ 移到全序 $\prec_i$ 中紧靠 $s$ 之前所得的全序；以 $\chi^u$ 表示元素 $u$ 的关联向量（$u\in U$）。计算 $0\le\epsilon\le-x(t)$，并将 $x':=x+\epsilon(\chi^t-\chi^s)$ 显式写成至多 $n$ 个向量的凸组合，所用向量选自 $b^{\prec_1},\ldots,b^{\prec_k}$ 及满足 $s\prec_i u\preceq_i t$ 的 $b^{\prec_i^{s,u}}$；另外要求若 $x'(t)<0$，则凸组合中不出现 $b^{\prec_i}$。
\item 令 $x:=x'$，将其凸组合中的向量重新命名为 $b^{\prec_1},\ldots,b^{\prec_k}$，返回步骤2。
\end{enumerate}
\st{定理14.20（Schrijver [2000]）} Schrijver 算法是正确的。
\pr 算法在 $D$ 中没有从 $P$ 到 $N$ 的路时终止，输出从 $P$ 不可达的顶点集 $X$。显然 $N\subseteq X\subseteq U\setminus P$，故对每个 $W\subseteq U$，$\sum_{u\in X}x(u)\le\sum_{u\in W}x(u)$。此外，没有边进入 $X$，所以或者 $X=\varnothing$，或者对每个 $j\in\{1,\ldots,k\}$ 都存在 $v\in X$，使 $X=\{u\in U:u\preceq_jv\}$。由定义，所有 $j$ 都满足 $\sum_{u\in X}b^{\prec_j}(u)=f(X)$。又由命题14.10，对所有 $W\subseteq U$、$j\in\{1,\ldots,k\}$，有 $\sum_{u\in W}b^{\prec_j}(u)\le f(W)$。因此
\begin{align*}
f(W)&\ge\sum_{j=1}^k\lambda_j\sum_{u\in W}b^{\prec_j}(u)
=\sum_{u\in W}\sum_{j=1}^k\lambda_jb^{\prec_j}(u)=\sum_{u\in W}x(u)\\
&\ge\sum_{u\in X}x(u)=\sum_{u\in X}\sum_{j=1}^k\lambda_jb^{\prec_j}(u)
=\sum_{j=1}^k\lambda_j\sum_{u\in X}b^{\prec_j}(u)=f(X),
\end{align*}

\srcpage{335}
这证明了 $X$ 为最优解。\qed
\st{引理14.21（Schrijver [2000]）} 每次迭代可在 $O(n^3+\gamma n^2)$ 时间内完成，其中 $\gamma$ 是一次预言机调用所需的时间。
\pr 只需证明步骤5可在该时间内完成。设 $x=\sum_{j=1}^k\lambda_jb^{\prec_j}$，且 $s\prec_i t$。先证明：
\textbf{断言：}对某个 $\delta\ge0$，$\delta(\chi^t-\chi^s)$ 可在 $O(\gamma n^2)$ 时间内写成向量 $b^{\prec_i^{s,v}}-b^{\prec_i}$（$s\prec_i v\preceq_i t$）的凸组合。

先作准备。设 $s\prec_i v\preceq_i t$。由定义，当 $u\prec_i s$ 或 $u\succ_i v$ 时，$b^{\prec_i^{s,v}}(u)=b^{\prec_i}(u)$。由次模性，当 $s\preceq_i u\prec_i v$ 时，
\begin{align*}
b^{\prec_i^{s,v}}(u)
&=f(\{w\in U:w\preceq_i^{s,v}u\})-f(\{w\in U:w\prec_i^{s,v}u\})\\
&\le f(\{w\in U:w\preceq_i u\})-f(\{w\in U:w\prec_i u\})=b^{\prec_i}(u).
\end{align*}
而对 $u=v$，
\begin{align*}
b^{\prec_i^{s,v}}(v)
&=f(\{w\in U:w\preceq_i^{s,v}v\})-f(\{w\in U:w\prec_i^{s,v}v\})\\
&=f(\{w\in U:w\prec_i s\}\cup\{v\})-f(\{w\in U:w\prec_i s\})\\
&\ge f(\{w\in U:w\preceq_i v\})-f(\{w\in U:w\prec_i v\})=b^{\prec_i}(v).
\end{align*}
还注意到 $\sum_{u\in U}b^{\prec_i^{s,v}}(u)=f(U)=\sum_{u\in U}b^{\prec_i}(u)$。
若某个 $s\prec_i v\preceq_i t$ 满足 $b^{\prec_i^{s,v}}=b^{\prec_i}$，断言显然成立；故可假设所有这样的 $v$ 都满足 $b^{\prec_i^{s,v}}(v)>b^{\prec_i}(v)$。递归定义
\[\kappa_v:=\frac{\chi^t_v-\sum_{v\prec_i w\preceq_i t}\kappa_w\bigl(b^{\prec_i^{s,w}}(v)-b^{\prec_i}(v)\bigr)}{b^{\prec_i^{s,v}}(v)-b^{\prec_i}(v)}\ge0
\qquad(s\prec_i v\preceq_i t).\]
可得 $\sum_{s\prec_i v\preceq_i t}\kappa_v(b^{\prec_i^{s,v}}-b^{\prec_i})=\chi^t-\chi^s$，因为对所有 $s\prec_i u\preceq_i t$，
\begin{align*}
\sum_{s\prec_i v\preceq_i t}\kappa_v\bigl(b^{\prec_i^{s,v}}(u)-b^{\prec_i}(u)\bigr)
&=\sum_{u\preceq_i v\preceq_i t}\kappa_v\bigl(b^{\prec_i^{s,v}}(u)-b^{\prec_i}(u)\bigr)\\
&=\chi^t_u,
\end{align*}
且全部分量之和为0。令 $\delta:=1/\sum_{s\prec_i v\preceq_i t}\kappa_v$，并将每个 $\kappa_u$ 乘以 $\delta$，断言得证。

现取 $\epsilon:=\min\{\lambda_i\delta,-x(t)\}$，$x':=x+\epsilon(\chi^t-\chi^s)$。若 $\epsilon=\lambda_i\delta\le-x(t)$，则
\[x'=\sum_{j=1}^k\lambda_jb^{\prec_j}+\lambda_i\sum_{s\prec_i v\preceq_i t}\kappa_v(b^{\prec_i^{s,v}}-b^{\prec_i}),\]
即已将 $x'$ 写成 $b^{\prec_j}$（$j\ne i$）及 $b^{\prec_i^{s,v}}$（$s\prec_i v\preceq_i t$）的凸组合。若 $\epsilon=-x(t)$，还可使用 $b^{\prec_i}$。最后按第4章习题5所示，在 $O(n^3)$ 时间内将凸组合缩减为至多 $n$ 个向量。\qed
\st{引理14.22（Vygen [2003]）} Schrijver 算法经过 $O(n^5)$ 次迭代后终止。
\srcpage{336}
\pr 若某次迭代的步骤5加入新向量 $b^{\prec_i^{s,v}}$ 后引入了边 $(v,w)$，则该次迭代中 $s\preceq_i w\prec_i v\preceq_i t$。所以 $d(w)\le d(s)+1=d(t)\le d(v)+1$，新边不会缩短从 $P$ 到任何 $v\in U$ 的距离。步骤5保证从不向 $P$ 添加元素，因此每个 $v\in U$ 的距离 $d(v)$ 从不减小。

将顶点对 $(t,s)$ 保持不变的一段连续迭代称为一个\textbf{块}。每个块包含 $O(n^2)$ 次迭代，因为其中每次迭代都使 $(\alpha,\beta)$ 按字典序下降。只需证明块数为 $O(n^3)$。

一个块只可能因以下至少一种原因结束（由 $t,s$ 的选择方式，且 $t=t^*$ 的迭代不会加入终点为 $t^*$ 的边，而顶点 $v$ 只有在 $v=s$ 时才能进入 $N$，此时 $d(v)<d(t)$）：
\begin{enumerate}[label=(\alph*)]
\item 某个 $v\in U$ 的距离 $d(v)$ 增加；
\item $t$ 从 $N$ 中删除；
\item $(s,t)$ 从 $A$ 中删除。
\end{enumerate}
分别计数。(a)型块显然有 $O(n^2)$ 个。

考虑(b)型。断言：对每个 $t^*\in U$，满足 $t=t^*$ 且 $x'(t)=0$ 的迭代有 $O(n^2)$ 次。这容易看出：任意两次这样的迭代之间，必有某个 $v\in U$ 的 $d(v)$ 改变；由于距离只能增加，这只能发生 $O(n^2)$ 次。因此(b)型阶段总数为 $O(n^3)$。

最后证明(c)型块有 $O(n^3)$ 个。只需说明：同一顶点对 $(s,t)$ 的下一个这种块出现之前，$d(t)$ 必改变。
对 $s,t\in U$，若 $(s,t)\notin A$ 或 $d(t)\le d(s)$，就称 $s$ 对 $t$ \textbf{无趣}（$t$-boring）。固定 $s^*,t^*\in U$，考虑从一个 $s=s^*,t=t^*$ 的块因 $(s^*,t^*)$ 被删去而结束之后，到下一次 $d(t^*)$ 改变之前的时间段。证明每个 $v\in\{s^*,\ldots,n\}$ 在此段始终对 $t^*$ 无趣；取 $v=s^*$ 即完成证明。

时间段开始时，由紧前一次迭代中 $s=s^*$ 的选择方式，每个 $v\in\{s^*+1,\ldots,n\}$ 都对 $t^*$ 无趣。由于 $(s^*,t^*)$ 被删除，$s^*$ 也对 $t^*$ 无趣。此段中 $d(t^*)$ 不变，而任意 $d(v)$ 从不减小，故只需检查新边的引入。
假设某次选取 $(s,t)$ 的迭代后，向 $A$ 加入边 $(v,t^*)$，其中 $v\in\{s^*,\ldots,n\}$。由证明开头的观察，该次迭代中 $s\preceq_i t^*\prec_i v\preceq_i t$，因此 $d(t^*)\le d(s)+1=d(t)\le d(v)+1$。分两种情形：若 $s>v$，则 $d(t^*)\le d(s)$，因为或者 $t^*=s$，或者 $s$ 对 $t^*$ 无趣且 $(s,t^*)\in A$。若 $s<v$，则 $d(t)\le d(v)$，因为或者 $t=v$，或者由 $s$ 的选择及 $(v,t)\in A$ 得到。两种情形都推出 $d(t^*)\le d(v)$，所以 $v$ 仍对 $t^*$ 无趣。\qed
定理14.20、引理14.21及引理14.22共同蕴含：
\st{定理14.23} 次模函数最小化问题可在 $O(n^8+\gamma n^7)$ 时间内求解，其中 $\gamma$ 为一次预言机调用的时间。\qed

\srcpage{337}
Iwata [2002] 描述了一个完全组合的算法（仅用加法、减法、比较和预言机调用，不用乘除法），并改进了运行时间（Iwata [2003]）。
\sect{14.5\quad 对称次模函数}
若对所有 $A\subseteq U$，$f(A)=f(U\setminus A)$，则称次模函数 $f:2^U\to\R$ \textbf{对称}。此特例的次模函数最小化问题是平凡的，因为对所有 $A\subseteq U$，
\[2f(\varnothing)=f(\varnothing)+f(U)\le f(A)+f(U\setminus A)=2f(A),\]
所以空集最优。因此只有排除这一平凡情形才有意义：寻找 $U$ 的非空真子集 $A$，使 $f(A)$ 最小。

Queyranne [1998] 推广第8.7节的算法，得到一个较简单的组合算法，仅用 $O(n^3)$ 次预言机调用。下述引理推广了引理8.38（习题14）。
\st{引理14.24} 给定对称次模函数 $f:2^U\to\R$，设 $n:=|U|\ge2$，可在 $O(n^2\theta)$ 时间内找到不同的 $x,y\in U$，使
\[f(\{x\})=\min\{f(X):x\in X\subseteq U\setminus\{y\}\},\]
其中 $\theta$ 是 $f$ 的预言机时间上界。
\pr 按以下方式对 $k=1,\ldots,n-1$ 构造次序 $U=\{u_1,\ldots,u_n\}$。假定已构造 $u_1,\ldots,u_{k-1}$，令 $U_{k-1}:=\{u_1,\ldots,u_{k-1}\}$。对 $C\subseteq U$，定义
\[w_k(C):=f(C)-\tfrac12\bigl(f(C\setminus U_{k-1})+f(C\cup U_{k-1})-f(U_{k-1})\bigr).\]
注意 $w_k$ 也对称。选取使 $w_k(\{u_k\})$ 最大的 $u_k\in U\setminus U_{k-1}$。最后令 $u_n$ 是 $U\setminus\{u_1,\ldots,u_{n-1}\}$ 中唯一的元素。显然次序的构造可在 $O(n^2\theta)$ 时间内完成。

\textbf{断言：}对所有 $k=1,\ldots,n-1$，以及不同的 $x,y\in U\setminus U_{k-1}$，若 $w_k(\{x\})\le w_k(\{y\})$，则
\[w_k(\{x\})=\min\{w_k(C):x\in C\subseteq U\setminus\{y\}\}.\]
对 $k$ 归纳。$k=1$ 时显然，因为对所有 $C\subseteq U$，$w_1(C)=\tfrac12f(\varnothing)$。
现设 $k>1$，且不同的 $x,y\in U\setminus U_{k-1}$ 满足 $w_k(\{x\})\le w_k(\{y\})$。另取 $Z\subseteq U$，使 $u_{k-1}\notin Z$，再取 $z\in Z\setminus U_{k-1}$。由 $u_{k-1}$ 的选取，$w_{k-1}(\{z\})\le w_{k-1}(\{u_{k-1}\})$，由归纳假设得 $w_{k-1}(\{z\})\le w_{k-1}(Z)$。此外，$f$ 的次模性给出：

\srcpage{338}
\begin{align*}
&(w_k(Z)-w_{k-1}(Z))-(w_k(\{z\})-w_{k-1}(\{z\}))\\
&=\tfrac12\bigl(f(Z\cup U_{k-2})-f(Z\cup U_{k-1})-f(U_{k-2})+f(U_{k-1})\bigr)\\
&\quad-\tfrac12\bigl(f(\{z\}\cup U_{k-2})-f(\{z\}\cup U_{k-1})-f(U_{k-2})+f(U_{k-1})\bigr)\\
&=\tfrac12\bigl(f(Z\cup U_{k-2})+f(\{z\}\cup U_{k-1})-f(Z\cup U_{k-1})-f(\{z\}\cup U_{k-2})\bigr)\\
&\ge0.
\end{align*}
故 $w_k(Z)-w_k(\{z\})\ge w_{k-1}(Z)-w_{k-1}(\{z\})\ge0$。
为完成断言的证明，取 $C\subseteq U$，满足 $x\in C$、$y\notin C$，分两种情形。

\textbf{情形1：}$u_{k-1}\notin C$。取上式中的 $Z=C$、$z=x$，得 $w_k(C)\ge w_k(\{x\})$。

\textbf{情形2：}$u_{k-1}\in C$。取 $Z=U\setminus C$、$z=y$，得 $w_k(C)=w_k(U\setminus C)\ge w_k(\{y\})\ge w_k(\{x\})$。

断言得证。将其用于 $k=n-1$、$x=u_n$、$y=u_{n-1}$，得
\[w_{n-1}(\{u_n\})=\min\{w_{n-1}(C):u_n\in C\subseteq U\setminus\{u_{n-1}\}\}.\]
对所有满足 $u_n\in C$、$u_{n-1}\notin C$ 的 $C\subseteq U$，有
\[w_{n-1}(C)=f(C)-\tfrac12\bigl(f(\{u_n\})+f(U\setminus\{u_{n-1}\})-f(U_{n-2})\bigr),\]
因此引理成立（取 $x:=u_n$、$y:=u_{n-1}$）。\qed
上述证明来自 Fujishige [1998]。现在可类似于定理8.39的证明继续。
\st{定理14.25（Queyranne [1998]）} 给定对称次模函数 $f:2^U\to\R$，可以在 $O(n^3\theta)$ 时间内找到使 $f(A)$ 最小的非空真子集 $A\subset U$，其中 $\theta$ 是 $f$ 的预言机时间上界。
\pr 若 $|U|=1$，问题平凡。否则应用引理14.24，在 $O(n^2\theta)$ 时间内找到不同的 $x,y\in U$，满足 $f(\{x\})=\min\{f(X):x\in X\subseteq U\setminus\{y\}\}$。接着递归寻找 $U\setminus\{x\}$ 的非空真子集，使 $f':2^{U\setminus\{x\}}\to\R$ 最小，其中 $y\notin X$ 时 $f'(X):=f(X)$，$y\in X$ 时 $f'(X):=f(X\cup\{x\})$。容易看出 $f'$ 对称且次模。

设非空真子集 $Y\subset U\setminus\{x\}$ 最小化 $f'$。由对称性，不妨设 $y\in Y$。断言 $\{x\}$ 或 $Y\cup\{x\}$ 最小化 $f$（在 $U$ 的全部非空真子集上）。任取满足 $x\in C$ 的 $C\subset U$。若 $y\notin C$，由 $x,y$ 的选择，$f(\{x\})\le f(C)$；若 $y\in C$，则 $f(C)=f'(C\setminus\{x\})\ge f'(Y)=f(Y\cup\{x\})$。因此对所有非空真子集 $C$，有 $f(C)\ge\min\{f(\{x\}),f(Y\cup\{x\})\}$。

为达到所述时间界，当然不能显式计算整个 $f'$。而是保存 $U$ 的一个划分，初始为全部单元素集。每层递归都把分别包含 $x,y$ 的两个块合并。这样便能借助 $f$ 的预言机高效计算 $f'$。\qed

\srcpage{339}
Nagamochi 和 Ibaraki [1998] 以及 Rizzi [2000] 进一步推广了该结果。
\sect{第14章习题}
\begin{enumerate}[label=\arabic*.]
\item 设 $G$ 是无向图，$M$ 是 $G$ 中的最大匹配。令 $\F$ 为满足以下条件的 $X\subseteq E(G)$ 的全体：存在关于 $M$ 的特殊花森林 $F$，使 $E(F)\setminus M=X$。证明 $(E(G)\setminus M,\F)$ 是贪心结构。\par
\textbf{提示：}使用第10章习题23。
\item 设 $(E,\F)$ 是贪心结构，$c':E\to\R_+$。对 $F\subseteq E$，考虑瓶颈函数 $c(F):=\min\{c'(e):e\in F\}$。证明贪心结构的贪心算法应用于 $(E,\F)$ 和 $c$ 时，能找到使 $c(F)$ 最大的 $F\in\F$。
\item 本题说明贪心结构也可定义为语言（参见定义15.1）。设 $E$ 是有限集。若字母表 $E$ 上的语言 $L$ 满足以下条件，则称其为\textbf{贪心结构语言}：
\begin{enumerate}[label=(\alph*)]
\item $L$ 包含空串；
\item 对所有 $(x_1,\ldots,x_n)\in L$ 及 $1\le i<j\le n$，$x_i\ne x_j$；
\item 对所有 $(x_1,\ldots,x_n)\in L$，$(x_1,\ldots,x_{n-1})\in L$；
\item 若 $(x_1,\ldots,x_n),(y_1,\ldots,y_m)\in L$ 且 $m<n$，则存在 $i\in\{1,\ldots,n\}$，使 $(y_1,\ldots,y_m,x_i)\in L$。
\end{enumerate}
若 $L$ 满足(a)、(b)、(c)及以下条件，则称为\textbf{反拟阵语言}：\par
(d$'$) 若 $(x_1,\ldots,x_n),(y_1,\ldots,y_m)\in L$，且 $\{x_1,\ldots,x_n\}\nsubseteq\{y_1,\ldots,y_m\}$，则存在 $i\in\{1,\ldots,n\}$，使 $(y_1,\ldots,y_m,x_i)\in L$。
\par 证明：字母表 $E$ 上的语言 $L$ 是贪心结构语言（反拟阵语言），当且仅当集合系统 $(E,\F)$ 是贪心结构（反拟阵），其中 $\F:=\{\{x_1,\ldots,x_n\}:(x_1,\ldots,x_n)\in L\}$。
\item 设 $U$ 有限，$f:2^U\to\R$。证明 $f$ 次模，当且仅当对所有 $X\subseteq U$、$y,z\in U$，有
\[f(X\cup\{y,z\})-f(X\cup\{y\})\le f(X\cup\{z\})-f(X).\]
\item 设 $P$ 是非空多拟阵。证明存在次模、单调函数 $f$，使 $f(\varnothing)=0$ 且 $P=P(f)$。（若所有 $A\subseteq B\subseteq E$ 都有 $f(A)\le f(B)$，则称 $f:2^E\to\R$ 单调。）
\item[$6.^*$] 证明非空紧集 $P\subseteq\R_+^n$ 是多拟阵，当且仅当：
\begin{enumerate}[label=(\alph*)]
\item 对所有 $0\le x\le y\in P$，有 $x\in P$；
\item 对所有 $x\in\R_+^n$，以及满足 $y,z\le x$、$y,z\in P$ 且对此性质极大的 $y,z$（即 $y\le w\le x$、$w\in P$ 蕴含 $w=y$，而 $z\le w\le x$、$w\in P$ 蕴含 $w=z$），有 $\one y=\one z$。
\end{enumerate}
\textbf{注：}这是 Edmonds [1970] 最初的定义。
\setcounter{enumi}{6}
\item 证明多拟阵贪心算法应用于向量 $c\in\R^E$ 和次模但未必单调的函数 $f:2^E\to\R$ 时，求得：
\end{enumerate}

\srcpage{340}
\[\max\left\{cx:\sum_{e\in A}x_e\le f(A)\text{ 对所有 }A\subseteq E\right\}.\]
\begin{enumerate}[label=\arabic*.,start=8]
\item 当 $f,g$ 是拟阵秩函数时，利用带权拟阵交算法产生的 $c_1,c_2$ 构造整数最优对偶解，从而证明定理14.12这一特例。（Frank [1981]）
\item[$9.^*$] 设 $S$ 是有限集，$f:2^S\to\R$。如下定义 $f':\R_+^S\to\R$：对任意 $x\in\R_+^S$，存在唯一的 $k\in\Z_+$、$\lambda_1,\ldots,\lambda_k>0$ 以及 $\varnothing\subset T_1\subset T_2\subset\cdots\subset T_k\subseteq S$，使 $x=\sum_{i=1}^k\lambda_i\chi^{T_i}$，其中 $\chi^{T_i}$ 是 $T_i$ 的关联向量。令 $f'(x):=\sum_{i=1}^k\lambda_i f(T_i)$。证明 $f$ 次模当且仅当 $f'$ 凸。（Lovász [1983]）
\setcounter{enumi}{9}
\item 设 $E$ 有限，$f:2^E\to\R_+$ 次模，且所有 $e\in E$ 满足 $f(\{e\})\le2$。（二元组 $(E,f)$ 有时称为\textbf{2-多拟阵}。）\textbf{多拟阵匹配问题}要求找出满足 $f(X)=2|X|$ 的最大基数集合 $X\subseteq E$（$f$ 当然由预言机给出）。
\par 设 $E_1,\ldots,E_k$ 是两两不交的无序对，$(E,\F)$ 是由独立性预言机给出的拟阵，其中 $E=E_1\cup\cdots\cup E_k$。\textbf{拟阵奇偶问题}要求找出最大基数的 $I\subseteq\{1,\ldots,k\}$，使 $\bigcup_{i\in I}E_i\in\F$。
\begin{enumerate}[label=(\alph*)]
\item 证明拟阵奇偶问题可多项式归约到多拟阵匹配问题。
\item[(b)$^*$] 证明多拟阵匹配问题可多项式归约到拟阵奇偶问题。\par
\textbf{提示：}使用次模函数最小化算法。
\item[(c)$^*$] 证明不存在运行时间关于 $|E|$ 为多项式的多拟阵匹配问题算法。（Jensen 和 Korte [1982]、Lovász [1981]）
\end{enumerate}
（一个问题多项式归约到另一个问题，意指前者可由使用后者预言机的多项式时间预言机算法求解，见第15章。）\par
\textbf{注：}Lovász [1980,1981] 给出了一个重要特例的多项式时间算法。
\item 若任意满足 $X\cap Y\ne\varnothing$、$X\cup Y\ne S$ 的 $X,Y\subseteq S$ 均有 $f(X)+f(Y)\ge f(X\cup Y)+f(X\cap Y)$，则称 $f:2^S\to\R\cup\{\infty\}$ 为\textbf{交叉次模函数}。
\par \textbf{次模流问题}如下：给定有向图 $G$，函数 $l:E(G)\to\R\cup\{-\infty\}$、$u:E(G)\to\R\cup\{\infty\}$、$c:E(G)\to\R$，以及交叉次模函数 $b:2^{V(G)}\to\R\cup\{\infty\}$。可行次模流是函数 $f:E(G)\to\R$，满足所有 $e\in E(G)$ 的 $l(e)\le f(e)\le u(e)$，以及
\[\sum_{e\in\delta^-(X)}f(e)-\sum_{e\in\delta^+(X)}f(e)\le b(X).\]
\end{enumerate}

\srcpage{341}
上式对所有 $X\subseteq V(G)$ 都须成立。任务是判定可行流是否存在；若存在，求使下列费用尽可能小的一个：
\[\sum_{e\in E(G)}c(e)f(e).\]
\par 证明此问题推广了最小费用流问题，以及在两个多拟阵之交上优化线性函数的问题。
\par \textbf{注：}Edmonds 和 Giles [1977] 引入的次模流问题可在强多项式时间内求解，见 Fujishige、Röck 和 Zimmermann [1989]；另见 Fleischer 和 Iwata [2000]。
\begin{enumerate}[label=\arabic*.,start=12]
\item[$12.^*$] 证明描述可行次模流的不等式系统（习题11）是 TDI。证明这蕴含定理14.12及19.10。（Edmonds 和 Giles [1977]）
\setcounter{enumi}{12}
\item 证明基多面体的顶点集合恰好是 $U$ 上所有全序 $\prec$ 所对应的向量 $b^\prec$ 的集合，其中
\[b^\prec(u):=f(\{v\in U:v\preceq u\})-f(\{v\in U:v\prec u\})\quad(u\in U).\]
\textbf{提示：}见定理14.11的证明。
\item 证明引理8.38是引理14.24的特例。
\end{enumerate}
\sect{第14章参考文献}
{\small\setlength{\parindent}{0pt}\setlength{\parskip}{5pt}\sloppy
\textbf{一般文献}\par
Bixby, R.E., and Cunningham, W.H. [1995]: Matroid optimization and algorithms. In: Handbook of Combinatorics; Vol. 1 (R.L. Graham, M. Grötschel, L. Lovász, eds.), Elsevier, Amsterdam, 1995\par
Björner, A., and Ziegler, G.M. [1992]: Introduction to greedoids. In: Matroid Applications (N. White, ed.), Cambridge University Press, Cambridge 1992\par
Fujishige, S. [1991]: Submodular Functions and Optimization. North-Holland, Amsterdam 1991\par
Korte, B., Lovász, L., and Schrader, R. [1991]: Greedoids. Springer, Berlin 1991\par
McCormick, S.T. [2004]: Submodular function minimization. In: Handbook on Discrete Optimization (K. Aardal, G. Nemhauser, R. Weismantel, eds.), Elsevier, Berlin (forthcoming)\par
Schrijver, A. [2003]: Combinatorial Optimization: Polyhedra and Efficiency. Springer, Berlin 2003, Chapters 44–49\par
\textbf{引用文献}\par
Edmonds, J. [1970]: Submodular functions, matroids and certain polyhedra. In: Combinatorial Structures and Their Applications; Proceedings of the Calgary International Conference on Combinatorial Structures and Their Applications 1969 (R. Guy, H. Hanani, N. Sauer, J. Schonheim, eds.), Gordon and Breach, New York 1970, pp. 69–87\par
Edmonds, J. [1979]: Matroid intersection. In: Discrete Optimization I; Annals of Discrete Mathematics 4 (P.L. Hammer, E.L. Johnson, B.H. Korte, eds.), North-Holland, Amsterdam 1979, pp. 39–49\par
Edmonds, J., and Giles, R. [1977]: A min-max relation for submodular functions on graphs. In: Studies in Integer Programming; Annals of Discrete Mathematics 1 (P.L. Hammer,\par
}

\srcpage{342}
{\small\setlength{\parindent}{0pt}\setlength{\parskip}{5pt}\sloppy
E.L. Johnson, B.H. Korte, G.L. Nemhauser, eds.), North-Holland, Amsterdam 1977, pp. 185–204\par
Fleischer, L., and Iwata, S. [2000]: Improved algorithms for submodular function minimization and submodular flow. Proceedings of the 32nd Annual ACM Symposium on the Theory of Computing (2000), 107–116\par
Frank, A. [1981]: A weighted matroid intersection algorithm. Journal of Algorithms 2 (1981), 328–336\par
Frank, A. [1982]: An algorithm for submodular functions on graphs. In: Bonn Workshop on Combinatorial Optimization; Annals of Discrete Mathematics 16 (A. Bachem, M. Grötschel, B. Korte, eds.), North-Holland, Amsterdam 1982, pp. 97–120\par
Fujishige, S. [1998]: Another simple proof of the validity of Nagamochi and Ibaraki’s min-cut algorithm and Queyranne’s extension to symmetric submodular function minimization. Journal of the Operations Research Society of Japan 41 (1998), 626–628\par
Fujishige, S., Röck, H., and Zimmermann, U. [1989]: A strongly polynomial algorithm for minimum cost submodular flow problems. Mathematics of Operations Research 14 (1989), 60–69\par
Grötschel, M., Lovász, L., and Schrijver, A. [1981]: The ellipsoid method and its consequences in combinatorial optimization. Combinatorica 1 (1981), 169–197\par
Grötschel, M., Lovász, L., and Schrijver, A. [1988]: Geometric Algorithms and Combinatorial Optimization. Springer, Berlin 1988\par
Iwata, S. [2002]: A fully combinatorial algorithm for submodular function minimization. Journal of Combinatorial Theory B 84 (2002), 203–212\par
Iwata, S. [2003]: A faster scaling algorithm for minimizing submodular functions. SIAM Journal on Computing 32 (2003), 833–840\par
Iwata, S., Fleischer, L., L., and Fujishige, S. [2001]: A combinatorial, strongly polynomialtime algorithm for minimizing submodular functions. Journal of the ACM 48 (2001), 761–777\par
Jensen, P.M., and Korte, B. [1982]: Complexity of matroid property algorithms. SIAM Journal on Computing 11 (1982), 184–190\par
Lovász, L. [1980]: Matroid matching and some applications. Journal of Combinatorial Theory B 28 (1980), 208–236\par
Lovász, L. [1981]: The matroid matching problem. In: Algebraic Methods in Graph Theory; Vol. II (L. Lovász, V.T. Sós, eds.), North-Holland, Amsterdam 1981, 495–517\par
Lovász, L. [1983]: Submodular functions and convexity. In: Mathematical Programming: The State of the Art – Bonn 1982 (A. Bachem, M. Grötschel, B. Korte, eds.), Springer, Berlin 1983\par
Nagamochi, H., and Ibaraki, T. [1998]: A note on minimizing submodular functions. Information Processing Letters 67 (1998), 239–244\par
Queyranne, M. [1998]: Minimizing symmetric submodular functions. Mathematical Programming B 82 (1998), 3–12\par
Rizzi, R. [2000]: On minimizing symmetric set functions. Combinatorica 20 (2000), 445– 450\par
Schrijver, A. [2000]: A combinatorial algorithm minimizing submodular functions in strongly polynomial time. Journal of Combinatorial Theory B 80 (2000), 346–355\par
Vygen, J. [2003]: A note on Schrijver’s submodular function minimization algorithm. Journal of Combinatorial Theory B 88 (2003), 399–402\par
}

